\documentclass[10pt,a4paper]{amsart}
\usepackage[margin=19mm]{geometry}
\usepackage{amsmath,amssymb,mathtools}
\usepackage[scr=rsfs,cal=euler]{mathalfa}
\usepackage{xfrac}
\usepackage[unicode,hidelinks,bookmarks=false]{hyperref}
\newcommand{\ie}{i.e.\ }
\newcommand{\cf}{cf.\ }
\newcommand{\resp}{respectively\ }
\newcommand{\Subsection}{\subsection}
\providecommand{\nobreakdash}{\nobreak\hbox{-}}
\setlength{\emergencystretch}{2em}
\multlinegap0pt
\usepackage{bm}
\usepackage{bbm}
\usepackage{dsfont}
\usepackage{xspace}
\usepackage{cancel}
\usepackage{mathtools}
\usepackage{amsthm}
\newtheorem{theorem}{Theorem}
\newtheorem{lemma}[theorem]{Lemma}
\title[Parametric decompositions into products of involutions]{Parametric decompositions into products of involutions}
\author{Gaofeng Huang, Frank Kutzschebauch, Tran Nam Son, Truong Huu Dung}
\date{Source: arXiv:2609.26808v1 (CC BY 4.0)}
\begin{document}
\maketitle

\noindent\emph{Source and licence.} Parametric decompositions into products of involutions. Version arXiv:2609.26808v1 (CC BY 4.0), distributed under the Creative Commons Attribution 4.0 International licence, as recorded in the arXiv licence metadata for this exact version and shown on the abstract page https://arxiv.org/abs/2609.26808v1. Full rights notice: licenses/arxiv-2609.26808-CC-BY-4.0.txt.\par\smallskip

\noindent\textit{Editorial scope.} Counted unit: the Lemma whose source label is \texttt{BiSo} in the article (source0.tex lines 299--307, source label \texttt{BiSo}), delivered together with the article's own complete original proof of it (source0.tex lines 309--384, one \texttt{proof} environment of 76 lines). No mathematics is written, repaired or re-derived: statement and proof are the article's own text line for line. Required context retained uncounted: none. Every definition, display and label of those lines is retained unchanged. Omitted from this package and not redistributed: 1--298, 308--308, 385--1060. Nothing inside a retained excerpt is omitted or altered: every retained body is delivered exactly as the article's lines are written, so no omission receipt of any kind accompanies this package. Citations: the retained excerpts carry no citation command at all, so no bibliography entry of the article is transcribed and no reference is redistributed. Reference adaptation: none was needed and none was made; every reference inside the delivered unit resolves inside it, so no unresolved reference or citation is produced. Printed slips kept literal: no printed word, symbol, hypothesis or number of the counted statement or of its proof was altered. Layout and class: the article's own class is \texttt{amsart} with packages including amsmath, amsxtra, amscd, amsthm, amsfonts, amssymb, mathdots, bbm, graphicx, longtable, booktabs, cite, tikz-cd, hyperref; the wrapper is the lane's pinned \texttt{amsart} editorial wrapper at 10pt on A4, which restates the theorem-like environments the retained text needs and adds the article's own macro definitions that the retained text needs (none), with extra wrapper packages (bm, bbm, dsfont, xspace, cancel, mathtools). The delivered numbering is the unit's own. Assets: no retained span contains a figure environment, a graphics-inclusion command, a PDF-page inclusion, a table or a TikZ picture; no image file is copied into this package, the package declares no asset dependency and no image file is redistributed with it. Licence: version arXiv:2609.26808v1 of this article is distributed under the Creative Commons Attribution 4.0 International licence, as recorded in the arXiv licence metadata for this exact version and shown on the abstract page https://arxiv.org/abs/2609.26808v1. A full rights notice is delivered at licenses/arxiv-2609.26808-CC-BY-4.0.txt. Characters: every retained excerpt is pure ASCII, so the delivered engine, XeLaTeX with its default Latin Modern text font, reports no missing character.
\begin{lemma}\label{BiSo}
	Let $R$ be a unital ring containing an infinite field $F$ which is central in $R$, and let $n>1$ be an integer.
	Suppose that $A\in M_n(R)$ is an upper triangular matrix whose diagonal entries
	$a_1,\ldots,a_n$ lie in $F$ and are pairwise distinct.
	Then there exists a matrix $P\in GL_n(R)$ such that
	\[
	P^{-1}AP=\mathrm{diag}(a_1,\ldots,a_n).
	\]
\end{lemma}

\begin{proof}
	We argue by induction on $n$.
	
	First, consider the case $n=2$. Let
	\[
	A=\begin{pmatrix}
		a_1 & y\\
		0 & a_2
	\end{pmatrix},
	\]
	where $y\in R$.
	Since $a_1,a_2\in F$ and $a_1\neq a_2$, the element $a_1-a_2$ is invertible in $F$.
	Set $x=y(a_1-a_2)^{-1}$.
	A direct computation shows that
	\[
	\begin{pmatrix}
		1 & x\\
		0 & 1
	\end{pmatrix}
	A
	\begin{pmatrix}
		1 & x\\
		0 & 1
	\end{pmatrix}^{-1}
	=
	\begin{pmatrix}
		a_1 & 0\\
		0 & a_2
	\end{pmatrix}.
	\]
	
	Now assume that $n>2$ and that the statement holds for all matrices of size less than $n$.
	Write
	\[
	A=\begin{pmatrix}
		a_1 & \alpha\\
		0 & A'
	\end{pmatrix},
	\]
	where $A'\in M_{n-1}(R)$ is an upper triangular matrix whose diagonal entries
	are $a_2,\ldots,a_n\in F$, and $\alpha$ is a row vector of length $n-1$.
	
	Let $A''=\mathrm{diag}(a_2,\ldots,a_n)$.
	By the induction hypothesis, there exists $P\in GL_{n-1}(R)$ such that
	\[
	P^{-1}A'P=A''.
	\]
	Since the elements $a_1-a_i$ are nonzero and belong to $F$ for all $i\ge2$,
	the diagonal matrix $a_1 I_{n-1}-A''$ is invertible in $M_{n-1}(R)$.
	
	Set
	\[
	x'=-\alpha P (a_1 I_{n-1}-A'')^{-1}
	\]
	and define
	\[
	P'=
	\begin{pmatrix}
		1 & 0\\
		0 & P
	\end{pmatrix}
	\begin{pmatrix}
		1 & x'\\
		0 & I_{n-1}
	\end{pmatrix}.
	\]
	A straightforward computation shows that
	\[
	P'^{-1}AP'=
	\begin{pmatrix}
		a_1 & 0\\
		0 & A''
	\end{pmatrix}.
	\]
	This completes the proof.
\end{proof}

\end{document}
